\clearpage
\begin{trapbox}
\textbf{C088 易错、边界与最小反例。}
\begin{center}% C088 PDF teaching figure; exact coordinates, equal x/y units.
\begin{tikzpicture}[x=0.78cm,y=0.78cm,>=stealth,font=\small]
\draw[rounded corners=5pt,black!25] (-4.2,-6.2) rectangle (12.8,5.9);
\node[font=\bfseries] at (4.3,5.25) {同尺度比较：同号与异号、竖直与水平都适用（$k=2/3$）};
\begin{scope}[shift={(0,0)}]
\node at (0,4.1) {同号：$x_1=x_2=3$};
\draw[->,black!45] (-3.6,0)--(3.7,0) node[right] {$x$};
\draw[->,black!45] (0,-2.8)--(0,2.75) node[above] {$y$};
\draw[blue!65!black,thick] (-3.3,-2.1999999999999997)--(3.3,2.1999999999999997);
\draw[violet!75!black,thick,dashed] (-3.3,2.1999999999999997)--(3.3,-2.1999999999999997);
\fill[teal!18] (0,0)--(3,2)--(3,-2)--cycle;
\draw[orange!85!black,very thick] (3,-2.6)--(3,2.6);
\fill (0,0) circle (1.7pt);\node at (.4,-.9) {$O$};
\fill (3,2) circle (1.7pt) node[above left=5pt] {$A$};
\fill (3,-2) circle (1.7pt) node[below left=5pt] {$B$};
\node at (0,-3.65) {竖直截线 $x=3$};
\node at (0,-5.2) {$S=\dfrac12\times4\times3=6$};
\end{scope}
\begin{scope}[shift={(8.6,0)}]
\node at (0,4.1) {异号：$x_1=3,x_2=-3$};
\draw[->,black!45] (-3.6,0)--(3.7,0) node[right] {$x$};
\draw[->,black!45] (0,-2.8)--(0,2.75) node[above] {$y$};
\draw[blue!65!black,thick] (-3.3,-2.1999999999999997)--(3.3,2.1999999999999997);
\draw[violet!75!black,thick,dashed] (-3.3,2.1999999999999997)--(3.3,-2.1999999999999997);
\fill[teal!18] (0,0)--(3,2)--(-3,2)--cycle;
\draw[orange!85!black,very thick] (-3.5,2)--(3.5,2);
\fill (0,0) circle (1.7pt);\node at (.4,-.9) {$O$};
\fill (3,2) circle (1.7pt) node[above left=5pt] {$A$};
\fill (-3,2) circle (1.7pt) node[above right=5pt] {$B$};
\node at (0,-3.65) {水平截线 $y=2$};
\node at (0,-5.2) {$S=\dfrac12\times6\times2=6$};
\end{scope}

\end{tikzpicture}
\end{center}
左图横坐标相同，但截点不同；右图两点都在 $x$ 轴上方，横坐标却异号。
两种情况都给正面积 $6$。

\begin{itemize}
\item \textbf{错把最终结果再乘 $\tfrac12$。}
面积行列式已经带 $\tfrac12$，而两条相反斜率使行列式出现 $2$。
最小反例：$k=1$、$A(1,1)$、$B(1,-1)$，真实面积为 $1$，而非 $\tfrac12$。

\item \textbf{漏掉绝对值。}
例如 $k=1$、$A(1,1)$、$B(-1,1)$，真实面积 $1$；
直接用 $kx_1x_2=-1$ 得到负数，说明只算了有向面积。

\item \textbf{误认 $A,B$ 一定一上一下或横坐标同号。}
正斜率渐近线经过第一、第三象限，负斜率渐近线经过第二、第四象限，
因此 $x_1x_2$ 可以小于 $0$；例2就是最小检验。

\item \textbf{误把渐近线截点当成双曲线上点。}
渐近线上的非原点点满足 $x^2/a^2-y^2/b^2=0$，
而双曲线满足 $x^2/a^2-y^2/b^2=1$，两者不能混为一谈。

\item \textbf{误认为 $x_1=x_2$ 会导致面积为零。}
当 $x_1=x_2=t\ne0$ 时，$A(t,kt)$、$B(t,-kt)$ 是\textbf{两个不同点}，
$\ell:x=t$，面积正好是 $kt^2$。

\item \textbf{误设截线必与 $x$ 轴相交。}
水平线 $y=n$（$n\ne0$）也有两个截点；
此时不能采用 $x$ 轴截距的推理，但核心面积公式仍适用。

\item \textbf{忽略真正的相交条件。}
若 $\ell$ 经过 $O$，交两条渐近线所得的有限截点都是 $O$，
不存在题意要求的两个不同截点；若 $\ell$ 与某渐近线重合，
交点不唯一，不能把它当成两条正常的截线交点。

\item \textbf{把 $b/a$ 写成 $a/b$。}
对 $H:x^2/a^2-y^2/b^2=1$，渐近线斜率绝对值是
$\boxed{b/a}$，不是 $a/b$；只有方程换成纵轴开口型时才需重算。
\end{itemize}

\medskip
\textbf{跨学段误区与最小反例：}
\begin{itemize}
\item \textbf{误认为必须先知道双曲线方程才能推出面积。}
不必。初中两条 $y=\pm kx$ 已能推出 $S=k|x_1x_2|$。
只在最后一步，双曲线才告诉我们 $k=b/a$。
\item \textbf{共底面积比忘了取绝对值。}
例如 $A=(2,2),\ B=(-1,1),\ C=(2,-2)$（$k=1$）。
此时 $B,C$ 在反向射线上，$OB/OC=1/2$，
面积比为 $1/2$，绝非 $-1/2$。
\item \textbf{把初中 $xy=q$ 的面积定值直接套给本题。}
初中的 $|q|/2$ 是单个\emph{反比例函数图象上的点}与坐标轴围成的直角三角形；
C088 的 $A,B$ 是渐近线截点，且任意截线面积并不定值。
\item \textbf{误以为坐标变换前后面积相等。}
在 $X=x/a+y/b,\ Y=x/a-y/b$ 下，$XY=1$，但
$S_{xy}=(ab/2)S_{XY}$；漏掉这个缩放因子会得出错误面积。
\item \textbf{把“切线面积定值”误当作“任意截线面积定值”。}
取 $a=3,b=2$，竖直线 $x=3$ 给面积 $6$，
此时 $x=3$ 本来就是顶点切线；换成割线 $x=6$，
面积就为 $(2/3)\cdot36=24\ne6$。
\emph{切线}是保证整族截线面积统一为 $ab$ 的一种充分条件，不能仅凭某一截线的面积等于 $ab$ 就判定它为切线。
\end{itemize}

\textbf{考试检查链：}
\[
\boxed{y_A=kx_1,\ y_B=-kx_2}
\quad\Longrightarrow\quad
\boxed{2S=|x_1y_2-x_2y_1|}
\quad\Longrightarrow\quad
\boxed{S=k|x_1x_2|}.
\]
\end{trapbox}
